\begin{abstract}
  离线强化学习从行为策略收集的静态数据集中学习策略，避免了在线交互的高昂成本与安全风险，在医疗诊断、自动驾驶、机器人控制等实际应用场景中具有重要价值。然而，由于行为策略通常是次优的，且离线数据集覆盖范围有限，分布偏移问题成为制约离线强化学习发展的核心挑战：当策略访问离线数据未覆盖的状态-动作对时，价值函数的估计将产生外推误差，进而导致策略性能退化。针对这一问题，本文以不确定性量化为核心手段展开系统研究：首先解决纯离线设定下的保守策略优化问题；在此基础上，针对离线数据覆盖边界对策略性能的根本性约束，进一步研究如何通过定向探索突破数据边界，形成从离线保守优化到在线高效微调的完整技术方案。

  在保守策略优化方面，本文提出基于自适应不确定性度量的离线强化学习算法。针对现有基于模型方法仅考虑状态转移不确定性、采用静态惩罚系数的局限，引入累积状态不确定性度量和自适应动态权重机制。累积状态不确定性通过递推公式综合考虑历史不确定性的累积效应，使智能体能够更全面地评估决策风险；自适应动态权重机制基于Sigmoid函数动态调节惩罚强度，实现探索与保守的自适应平衡。理论分析证明累积不确定性的单调性和自适应权重的优越性。实验结果表明，该方法在D4RL MuJoCo基准测试集上显著优于MOPO等基线方法。

  上述纯离线方法的策略性能仍受限于离线数据的覆盖范围。为突破这一根本性约束，本文进一步提出基于定向探索模型的离线到在线强化学习算法。该方法将不确定性信息从“惩罚信号”转变为“探索引导信号”：在离线阶段通过集成模型进行保守策略优化，在在线阶段显式引导策略探索高不确定性且高价值的区域。提出双重不确定性估计方法区分偶然不确定性和认知不确定性，设计动态权重自适应机制平滑过渡探索与利用。理论分析证明定向探索策略相比均匀采样具有更优的收敛速率。实验结果表明，该方法仅需50k步在线交互即可显著优于现有基线方法。

  \keywords{离线强化学习，不确定性估计，分布偏移，定向探索，离线到在线强化学习}

\end{abstract}

















\begin{englishabstract}

  Offline reinforcement learning enables policy learning from pre-collected static datasets, avoiding the high costs and safety risks of online interactions. This approach holds significant practical value in domains such as medical diagnosis, autonomous driving, and robotic control. However, since behavioral strategies are generally suboptimal and offline datasets have limited coverage, the core challenge hindering offline reinforcement learning: when the learned policy visits state-action pairs not covered by offline data, value function estimation produces extrapolation errors, leading to performance degradation. To address this problem, this thesis systematically studies uncertainty quantification: first solving the conservative policy optimization problem in purely offline settings; then, to overcome the fundamental constraint that offline data boundaries impose on policy performance, investigating how directed exploration can transcend these boundaries, forming a complete technical solution from offline conservative optimization to efficient online fine-tuning.

  In conservative policy optimization, an offline reinforcement learning algorithm based on adaptive uncertainty measurement is proposed. To address the limitations of existing model-based methods that only consider state transition uncertainty and employ static penalty coefficients, cumulative state uncertainty measurement and an adaptive dynamic weight mechanism are introduced. The cumulative state uncertainty comprehensively accounts for accumulated historical uncertainty through a recursive formula, enabling more thorough risk assessment. The adaptive dynamic weight mechanism based on Sigmoid function dynamically adjusts penalty strength to achieve better exploration-conservatism balance. Theoretical analysis proves the monotonicity of cumulative uncertainty and the superiority of adaptive weights. Experiments demonstrate significant improvements over MOPO on the D4RL MuJoCo benchmark.

  The above purely offline method's policy performance remains constrained by offline data coverage. To break through this fundamental limitation, an offline-to-online reinforcement learning algorithm based on directed exploration is further proposed. This method transforms uncertainty information from a "penalty signal" to an "exploration guide": during offline learning, conservative policy optimization is performed via ensemble models; during online learning, the policy is explicitly guided to explore high-uncertainty, high-value regions. A dual uncertainty estimation method distinguishes aleatoric and epistemic uncertainty, with a dynamic weight adaptation mechanism smoothing the transition between exploration and exploitation. Theoretical analysis proves superior convergence rates compared to uniform sampling. Experiments show significant improvements over baselines with only 50k online interaction steps.

  \englishkeywords{offline reinforcement learning, uncertainty estimation, distribution shift, directed exploration, offline-to-online reinforcement learning}

\end{englishabstract}
